\section{Source audit, verification, and integration status}
\label{ag:sec:audit}

\subsection{What was inspected and what was not}
The GitHub metadata anchor used for the final source inventory is
\[
 \texttt{\pin}.
\]
The relevant canonical chapters, the immediate predecessor's Gamma
sector, and its research-question section were read as TeX through the
repository connection. The source manifest in the accompanying package
records exact paths and the returned Git blob identifiers. Some initial
reads used the default branch; the predecessor's research questions and
Gamma-section identity were subsequently checked at the displayed
anchor. A Git blob identifier is not being presented as a SHA-256 hash
of a locally downloaded file.

The current incoming directory contained five ZIP archives, concerning
alternating tails, finite-difference determinants, higher quadratic jets,
higher spectral defects, and quartic spectral reconstruction. The
repository connection exposed their metadata but not their binary
contents. The ZIPs were \emph{not} extracted in this continuation.
Matching-title Library TeX abstracts were retrieved for the
finite-difference, higher-quadratic-jet, and higher-spectral-defect
reports. Their identity with the repository ZIP bytes was not
established. These abstracts concern determinantal or spectral
factorization problems; they are contextual comparisons, not audited
proof dependencies. No assertion that every incoming result has been
read, or that the entire repository has been certified, is made.

The main proof dependencies are therefore unusually explicit: the
functional-monomial decomposition, the Gamma Bell exponential, elementary
weighted integration by parts, and the standard Hurwitz and Lerch
identities. The necessary algebraic decomposition is reproved here.
The one concrete predecessor question used to measure progress is
question~5 of \cite[Section~8]{HTG}, concerning the next Gamma
projection. Repository files have not been modified remotely.

\subsection{Classical inputs and the boundary of the novelty claim}
The functional-monomial construction is an established component of the
integro-differential-algebra literature \cite{GRR}; its adoption in the
canonical manuscript and predecessor is retained with attribution
\cite{Manuscript,HTG}. The pure derivative-square coefficients, the
low-index function $f_0$, and the entire $r=2$ family are inherited
specializations. They should not be relabeled as new theorems during
integration.

Polynomially weighted Hurwitz integrals and their relations to
Bernoulli polynomials, log Gamma, and polygamma functions have extensive
classical antecedents; the papers of Espinosa and Moll explicitly
study definite and indefinite integrals of this kind \cite{EM1,EM2}.
The differentiation, translation, Mellin, Bernoulli, and functional-
equation identities used here are standard \cite{DLMFzeta,DLMFpsi,
DLMFlerch,DLMFBernoulli,DLMFintegers}. Repeated integration by parts
itself is not a new method.

The proposed additions relative to the inspected proof chain are the
all-$r$ affine square-sector coefficient theorem, the arbitrary symmetric
profile extension, their common beta--Hurwitz transport, and the
explicit selector classification as a statement over $\C(s)$. The
all-order Stieltjes coefficient rules and anchored primitive formulas
are proved consequences in the stated normalization. Particular
low-order integral values may already occur in the literature.
This report does not establish global historical priority for each
formula, nor does it claim a new special-function algebraic-independence
result at fixed numerical arguments.

\subsection{The exact proof-status distinctions}
The following status list should accompany any manuscript integration.

\paragraph{Proved here.}
Theorem~\ref{ag:thm:main} gives every
$u_0^a u_1^k u_r^2$ coefficient in the canonical inverse-Gamma normal
form. Section~\ref{ag:sec:profiles} gives its general symmetric-profile
analogue. The Hurwitz and Lerch transforms, their removable-pole
prescriptions, all Stieltjes orders, and all anchored shift primitives
are analytic identities on their specified domains. Theorem~
\ref{ag:thm:selector-classification} classifies exactly the symmetric,
endpoint-vanishing polynomial weights with the specified single-zeta
functional closure. Its necessity uses a proof of independence over
$\C(s)$, rather than independence of a declared formal list.

\paragraph{A source question partially resolved.}
The next-Gamma-projection question is answered for $u_3^2$ with arbitrary
powers of $u_0,u_1$ and no additional $u_2$ factor, and simultaneously
for every higher derivative square. Terms such as
$u_0^a u_1^k u_2^b u_3^2$ with $b>0$ lie beyond the quadratic
high-derivative sector used here. The complete $\sigma_3$ projection is
not claimed.

\paragraph{Not resolved by this report.}
The unrestricted unequal-even-tail optimal-prefix conjecture, the
Gaussian $S_6$ target, and the current revised $S_8$ candidate remain
separate. The formal Gamma theorem does not evaluate all of its
finite-part moments. No proof-assistant formalization, independent
human peer review, or certified numerical enclosure is claimed.

\subsection{Corrections, safeguards, and editorial observations}
No false inherited mathematical theorem was established by the targeted
audit. Several errors would arise from tempting but invalid extensions,
and the package supplies explicit guards against them.
First, normal reduction must precede truncation of higher derivatives:
$\NF(u_0u_4)=u_2^2$, whereas truncating $u_4$ beforehand gives zero.
Second, $\NF$ is not multiplicative. Third, weighted symbol reduction
requires the Appell primitive property in
Lemma~\ref{ag:lem:appell-primitive}; allowing an arbitrary
$x$-dependent primitive would destroy the quotient argument. Fourth,
removable poles in endpoint formulas must be canceled jointly, not
by assigning independent finite parts to their summands. Finally,
$B_1(1)=+1/2$ and $B_1(0)=-1/2$ play different roles in the selector
calculation. These are proof obligations and implementation safeguards,
not accusations that the predecessor omitted them.

A minor editorial duplication was observed in the inspected manuscript
README: the paragraph describing the full-Cayley, 1,032-input first-$X$
family appeared twice. Removing one copy would improve the overview
without changing any mathematical claim. That suggestion should be
checked against the current README before editing, since the repository
is changing rapidly. No correction to a numbered theorem is proposed
on the basis of that duplication.

\subsection{Exact finite verification and numerical diagnostics}
The executable checks distinguish four kinds of exact finite assertion.
\begin{center}
\small
\begin{tabular}{@{}p{0.25\linewidth}rp{0.57\linewidth}@{}}
\toprule
Replay & Assertions & Scope\\
\midrule
\src{check_algebra.py}&1,255&Normal forms and primitives through weight 18;
independent Bell sectors; all-index formula specializations; kernel
coefficient identities; deliberately wrong alternatives rejected.\\[4pt]
\src{check_endpoint_algebra.py}&728&Beta endpoint derivatives and residues;
complete-harmonic coefficients; Bernoulli endpoint conventions and
rational offsets.\\[4pt]
\src{check_profiles.py}&75&Three independent polynomial profiles compared
with direct Bell expansion and canonical reduction.\\[4pt]
\src{check_selectors.py}&129&Selector basis, beta--Bernoulli transport,
Stieltjes residues, inverse expansions, and an endpoint-sign guard.\\
\midrule
Total&2,187&Finite exact rational or symbolic assertions.\\
\bottomrule
\end{tabular}
\end{center}
The algebra replay compares both sides with independently constructed
representations. In particular it checks $A_m=N_m+\D T_m$, not merely
the displayed selected coefficient. The profile replay does not invoke
the proposed profile formula to create its Bell-polynomial reference.
The infinite-family assertions nevertheless rest on the mathematical
proofs in the article, not the largest computed index.

The numerical program has four separately runnable parts: 15 kernel,
14 transform, 6 Stieltjes, and 21 Lerch/selector comparisons, totaling
56. At 50 working decimal digits it compares beta integrals with
hypergeometric values, compact Hurwitz averages with finite endpoint
expressions, Mellin integrals with both, and polygamma/Stieltjes
quadratures with their explicit evaluations. It includes a complex
spectral argument and rational-endpoint root-of-unity filters.
Errors are scaled by $\max(1,|L|,|R|)$; small absolute values therefore
do not imply the same number of relative digits.

The largest recorded scaled residual is approximately
$2.62\times10^{-29}$ in a comparison of two \emph{different quadrature
orders} for the first Stieltjes average. Its declared tolerance is
$10^{-23}$. The comparison of the finer quadrature with the finite
endpoint formula is about $5.49\times10^{-43}$, and the explicit
quartic $\gamma_1$ evaluation agrees to about $4.22\times10^{-40}$
in the same scaled measure. The default tolerance for the other
comparisons is $10^{-30}$, with explicit local overrides documented
in the program. These are floating-point diagnostics, not interval
error bounds.

All scripts, exact coefficient tables, normal-form primitives,
individual results, environment information, and a replay entry point
are included. The integration fragment preserves unique \src{ag:}
labels, the distinction between ordinary and finite-part integrals,
and the nonmultiplicative meaning of the normal projection.
